\section{Discussion}
\label{sec:discussion}

\para{Why grounding reshapes rather than improves.} A single qualitative case captures the mechanism. On the Factuality topic, from the same seed, the \smallexp scaffold produced ``Chain-of-Verification''---a real, already-published method---whereas \litcomp and \deduction produced more elaborated ``counterfactual chain-of-reasoning'' and ``counterfactual retrieval'' variants. Asking the model to predict whether a small test would pass optimizes for ideas that \emph{will} survive scrutiny, \ie safe, known techniques: lower novelty, higher feasibility. Asking it to flip an assumption or out-distance retrieved papers optimizes for differentiation: higher judged novelty, lower feasibility. Grounding is therefore best understood as a \emph{steering knob on the novelty--feasibility axis}, not a uniform quality lever. This reframes the survey ranking (graph-grounded $>$ Chain-of-Ideas $>$ vanilla): the apparent advantage of grounding may be a judged-novelty effect that our objective metric does not corroborate.

\para{The mirage is the most important result.} Judged novelty rises for \deduction and \litcomp, yet objective distance to prior work does not, and the two measures are uncorrelated at the item level ($\rho = 0.01$). Because the two off-family judges \emph{agree with each other} ($\rho = 0.41$), the mirage is not judge noise---they reliably reward something real, just not occupancy of unexplored regions of the literature. The reward goes to \emph{rhetorical} novelty: framing, elaborate mechanism descriptions, and signal words like ``counterfactual.'' This operationalizes the \emph{All That Glitters} warning~\citep{glitters2025} and reproduces RQ-Bench's mirage~\citep{rqbench2026} in a controlled setting where the idea set is held fixed and only the measurement changes. The practical consequence is direct: measure novelty as distance to retrieved prior work, not as an LLM's opinion, and never report a single-judge novelty number alone.

\para{Calibration against prior work.} Our control means (judge novelty $\approx 5.5$, feasibility $\approx 7.2$) sit in a plausible band relative to the human AI-condition means of \citet{si2024can} (novelty $\approx 5.6$, feasibility $\approx 6.3$), so the judges are not wildly miscalibrated in absolute terms---the problem is what they track, not their scale.

\para{Failure modes.} We observed three. (1)~\emph{Topic-anchored convergence}: ``counterfactual'' dominated the Factuality topic across conditions, with grounding modulating elaboration around the motif rather than escaping it. (2)~\emph{Rediscovery}: \smallexp repeatedly produced named existing methods. (3)~\emph{Scale compression}: the absolute rubric compresses ideas into the 5--6 band, hiding preferences that pairwise judging readily reveals (\eg \litcomp's flat absolute overall score but $79\%$ win-rate).

\para{Limitations.} (1)~\emph{Single backbone.} We fixed \gptfourone to isolate grounding, trading external validity for internal validity; mechanism effects may differ for other model families. (2)~\emph{No item-level human anchor.} The released Si review data is anonymized without idea text, so we could not correlate our judges against human scores item-by-item; we validated via inter-judge agreement and convergent validity instead, but the absent human anchor is a real gap. (3)~\emph{Thought-experiment grounding.} Generic ideas cannot be cheaply executed in-loop, so \smallexp is predict-and-critique rather than a real run; a domain where ideas \emph{are} runnable (\eg a fixed prompting benchmark) could test true experimental grounding. (4)~\emph{\lgn is one embedding model's view}; cosine distance to abstracts is a proxy, not ground truth---but it is judge-independent, so its disagreement with the judge is informative regardless of absolute calibration. (5)~\emph{Scale and scope.} $N = 5$ per cell gives medium power; under-powered trends (\eg \outcome novelty, $p = 0.09$) are suggestive only, and we tested single-mechanism scaffolds, not the 2--3-mechanism combinations most strong systems use. (6)~\emph{Contamination.} The public topics are within \gptfourone's training data, which may inflate the rediscovery failure mode.

\para{Broader implications.} Grounding is a cheap, deployable prompting protocol, and our results let practitioners pick the knob that matches their goal rather than chasing a single ``quality'' number. More broadly, the mirage result is a caution for the rapidly growing practice of using LLM judges to score machine-generated research: a metric that two strong judges agree on can still be measuring the wrong construct.
